\section{Compatibility of the Lorentz action across field levels}
\label{sec:label-descent}
The continuation argument uses one Lorentz action on the union of the field
spaces. Constructing an action separately at every finite level would not
be enough: the changes of basis used to compare with the free theory need
not commute with the inclusions. The compatibility is instead proved from
the Euclidean action and the infinitesimal boost relation.

Write $E_d$ for the interacting finite-dimensional complex field space and
$\iota_{de}:E_d\to E_e$, $d\leq e$, for its injective inclusions. Let
$E_d^0$ be the corresponding reference spaces. The following input isolates
the finite-level comparison produced by the calibration argument.
\begin{assumption}[Calibrated finite-level actions]\label{input:calibrated-actions}
For each $d$ there is a complex-linear isomorphism $Q_d:E_d^0\to E_d$.
The reference spaces carry continuous representations $R_d^0$ of $O(4)$
and $T_d^0$ of the proper orthochronous Lorentz group. For each spatial
coordinate $a$, their generators satisfy
\[
 \left.\frac{d}{dt}T_d^0(B_a(t))\right|_{t=0}=-iL_{d,a}^0,
\]
where $L_{d,a}^0$ is the generator of the Euclidean time--$a$ plane
rotation, with the conventions of Section~\ref{sec:lorentz-covariance}.
On spatial rotations the two reference actions agree. Define
$R_d(g)=Q_dR_d^0(g)Q_d^{-1}$ and let $L_{d,a}$ be its indicated
Euclidean generator. These actions satisfy
\[
 R_e(g)\iota_{de}=\iota_{de}R_d(g),\qquad
 L_{d,a}Q_d=Q_dL_{d,a}^0.
\]
There are antilinear involutions $J_d,J_d^0$, compatible with inclusions,
such that $J_dQ_d=Q_dJ_d^0$ and $J_d^0$ commutes with $T_d^0$.
The inclusions form a directed system and the field union is its algebraic
direct limit.
\end{assumption}
The comparison maps $Q_d$ are the same maps that define $R_d$; a fresh
choice at the Lorentz stage would not justify the displayed generator
identity. Production of these maps and the reference actions is separate
from the compatibility proof below.

\begin{proposition}\label{prop:lorentz-label-descent}
Under Input~\ref{input:calibrated-actions}, the finite-level representations
\[
 T_d(\Lambda)=Q_dT_d^0(\Lambda)Q_d^{-1}
\]
commute with all inclusions. They therefore induce one representation
$T$ on the field union, continuous on each finite level. Its spatial
action is the Euclidean one, its boost generator on $E_d$ is
$-iL_{d,a}$, and it commutes with the induced involution $J$.
\end{proposition}
\begin{proof}
Conjugation immediately gives the representation laws and continuity on
each level. Differentiation at zero, using the fixed linear map $Q_d$,
gives the boost generator $K_{d,a}=-iL_{d,a}$. Differentiating the
Euclidean inclusion law along its plane rotation gives
\[
 L_{e,a}\iota_{de}=\iota_{de}L_{d,a},
 \qquad K_{e,a}\iota_{de}=\iota_{de}K_{d,a}.
\]
For $v\in E_d$, the two curves
\[
 f(t)=T_e(B_a(t))\iota_{de}v,\qquad
 g(t)=\iota_{de}T_d(B_a(t))v
\]
have the same initial value and solve $h'=K_{e,a}h$. The derivative at
arbitrary $t$ follows from the group law by translating the parameter in
the derivative at zero; for $g$ use the generator intertwining just proved.
Uniqueness of this constant-coefficient equation can be seen directly:
the derivative of $e^{-tK_{e,a}}(f(t)-g(t))$ is zero, and its value at
zero is zero. Hence the inclusion commutes with every coordinate boost.

It commutes with each spatial rotation because the transported spatial
action is $R_d$ and the Euclidean inclusion law already holds.
For fixed $d,e$, the transformations commuting with $\iota_{de}$ form
a subgroup. Products follow by two substitutions. For inverses, if
$T_e(\Lambda)\iota_{de}=\iota_{de}T_d(\Lambda)$, multiply on the left
and right by the respective inverse transformations. Spatial rotations
and coordinate boosts generate the full proper orthochronous group, by
the explicit reduction of $\Lambda e_0$ in the proof of
Theorem~\ref{thm:lorentz-covariance}. Thus every transformation commutes
with the inclusions. Notice that this argument has not asserted any
compatibility of the maps $Q_d$ themselves.

Define the action on a union element represented by $v\in E_d$ to be
represented by $T_d(\Lambda)v$. If another representative lies in $E_e$,
pass to a common higher level; the inclusion compatibility proves the
same resulting class. This proves well-definedness. Linearity and the
group laws hold after putting finitely many representatives in a common
level, and the inverse action is that of $\Lambda^{-1}$. The asserted
spatial and derivative identities follow in the same way. Finally
\[
 J_dT_d(\Lambda)Q_dv
 =Q_dJ_d^0T_d^0(\Lambda)v
 =Q_dT_d^0(\Lambda)J_d^0v
 =T_d(\Lambda)J_dQ_dv.
\]
Surjectivity of $Q_d$ and compatibility of the involutions prove
commutation with $J$ on the union.
\end{proof}
